\section{Introduction}

Many contemporary distributed systems operate in environments where data and computational resources are owned by different organizations. In such settings, data exchange must adhere to legal agreements, privacy regulations, and organizational policies. These constraints are common in domains such as healthcare, scientific research, and public-sector infrastructures, where unrestricted data exchange is not acceptable.

DYNAMOS (Dynamically Adaptive Microservice-based Operating System) is a data exchange middleware developed by the Complex Cyber-Infrastructure (CCI) group at the University of Amsterdam. The system implements a set of recurring data sharing archetypes that describe how data and computation are distributed among participants under policy constraints. DYNAMOS realizes these archetypes by dynamically composing chains of microservices that execute data sharing requests in compliance with formally evaluated policies \cite{stutterheim2023thesis,stutterheim2023icsa}. The middleware is designed to be adaptive, allowing architectural and execution decisions to change at runtime in response to policy requirements.

The current implementation of DYNAMOS uses gRPC for inter-service communication and relies primarily on atomic request-response interactions \cite{stutterheim2023thesis}. This communication model is suitable for small or bounded data transfers but becomes limiting when applied to large-scale or continuous data exchange. In practice, data in distributed systems is often produced and consumed incrementally, for example in sensor data streams, distributed analytics pipelines, or staged processing of large research datasets.

When request-response communication is used in such scenarios, systems must either buffer large volumes of data before transmission or split transfers into multiple independent requests. Buffering increases memory usage and can lead to idle time between producers and consumers. Repeated request-response interactions introduce additional coordination overhead. Both effects negatively impact scalability and resource utilization as data volumes and the number of participating services increase.

Streaming-based communication has been proposed as a solution to these limitations. Streaming allows data to be transferred and processed incrementally, enabling overlap between computation and communication and reducing peak memory usage. gRPC supports several streaming interaction patterns, including client-side, server-side, and bidirectional streaming.

Integrating streaming communication into DYNAMOS is not straightforward. Streaming relies on long-lived connections and continuous data flows, which must be reconciled with DYNAMOS' policy-driven execution model, dynamic microservice composition, orchestration decisions, and fault handling mechanisms. Without careful design, the introduction of streaming may interfere with policy enforcement or reduce the ability to evaluate and reason about system behavior.

The goal of this thesis is to investigate how streaming communication can be incorporated into DYNAMOS in a way that improves scalability and performance while preserving policy compliance and adaptive behavior. This involves analyzing existing streaming communication mechanisms, designing architectural extensions to the middleware, implementing streaming support, and evaluating the resulting system under realistic workloads.

The main research question addressed in this thesis is:

\begin{quote}
\textbf{Main RQ:} How can streaming communication be integrated into a dynamically adaptive, policy-driven microservice-based data exchange middleware such as DYNAMOS, while preserving compliance guarantees and improving scalability and performance?
\end{quote}

To answer this question, the following sub-research questions are defined:

\begin{itemize}
  \item \textbf{RQ1:}
  Which streaming communication technologies, frameworks, and interaction patterns are suitable for large-scale and continuous data exchange in microservice-based systems such as DYNAMOS?

  \item \textbf{RQ2:}
    How can streaming communication be integrated into a microservice-based architecture in a way that preserves dynamic service composition, policy-driven orchestration, and isolation guarantees?

  \item \textbf{RQ3:}
    What is the impact of the implemented streaming data paths on scalability, performance, and robustness compared to a traditional request-response approach under realistic DYNAMOS workloads?
\end{itemize}

